\documentclass[11pt]{article}
\usepackage{rd_audit_style}

\RDDocumentTitle{R\&D Experiment Notebook}
\RDCompany{<< document.company >>}
\RDFY{<< document.financial_year >>}
\RDPreparedBy{<< document.prepared_by >>}
\RDVersion{<< document.version >>}
\RDConfidentiality{<< document.confidentiality >>}

\begin{document}
\RDTitlePage
\tableofcontents
\clearpage

\section{How to use this notebook}
\RDSectionIntro{Create one section per independent run, test cycle, prototype iteration, benchmark round, or other discrete investigative step. The objective is to preserve the chronology of technical work and the logic of what was learned.}

\RDAuditNote{Do not rewrite history after the fact. Where a hypothesis changes, record the original assumption, why it changed, and what new experiment followed. A notebook should show the actual progression of technical investigation.}

\section{Notebook index}
\rowcolors{2}{white}{RDSoft}
\begin{xltabular}{\textwidth}{L{2.4cm}L{2.3cm}L{4.4cm}L{2.8cm}Y}
\rowcolor{RDBand}
\textbf{Run ID} & \textbf{Date} & \textbf{Experiment title} & \textbf{Outcome} & \textbf{Related project / uncertainty} \\
RUN-001 & {[Date]} & {[Title]} & {Confirmed / Rejected / Mixed / Ongoing} & {[Project / UT-ID]} \\
RUN-002 & {[Date]} & {[Title]} & {Confirmed / Rejected / Mixed / Ongoing} & {[Project / UT-ID]} \\
\end{xltabular}
\rowcolors{2}{}{}

\section{Audit-grade experiment run template}
\subsection{Administrative details}
Run ID: {[RUN-XXX]}\\
Date or date range: {[Date range]}\\
Engineer(s): {[Name(s)]}\\
Project reference: {[Project ID / Name]}\\
Uncertainty reference: {[UT-ID]}\\
Supporting issue / ticket: {[ID]}\\
Repo / branch / environment: {[Reference]}

\subsection{Technical objective}
State what this run was intended to determine.

\subsection{Background}
\RaggedRight
Summarise the prior system state, constraints, assumptions, and why this run was necessary.\par
\justifying

\subsection{Hypothesis}
Write the testable expectation in direct technical terms.

\subsection{Variables and conditions}
\begin{RDChecklist}
\item Inputs or datasets used.
\item Environment or infrastructure configuration.
\item Control condition or baseline.
\item Load, concurrency, payload, or data-volume assumptions.
\item Acceptance thresholds and failure triggers.
\end{RDChecklist}

\subsection{Method}
Describe the exact method used. Include the implementation approach, instrumentation, build version, and any benchmarking or observation steps.

\subsection{Execution log}
\rowcolors{2}{white}{RDSoft}
\begin{xltabular}{\textwidth}{L{2.4cm}L{2.4cm}Y}
\rowcolor{RDBand}
\textbf{Timestamp / step} & \textbf{Actor / tool} & \textbf{Action / observation} \\
{[Time]} & {[Engineer / script / service]} & {[What happened]} \\
{[Time]} & {[Engineer / script / service]} & {[What happened]} \\
\end{xltabular}
\rowcolors{2}{}{}

\subsection{Results}
\rowcolors{2}{white}{RDSoft}
\begin{xltabular}{\textwidth}{L{3.8cm}C{2.6cm}C{2.6cm}Y}
\rowcolor{RDBand}
\textbf{Metric} & \textbf{Baseline} & \textbf{Experiment} & \textbf{Notes} \\
Latency & {[value]} & {[value]} & {[Observation]} \\
Throughput & {[value]} & {[value]} & {[Observation]} \\
Error rate & {[value]} & {[value]} & {[Observation]} \\
Resource usage & {[value]} & {[value]} & {[Observation]} \\
\end{xltabular}
\rowcolors{2}{}{}

\subsection{Unexpected behaviour or failure modes}
Record any non-obvious outcomes, regressions, data anomalies, or side effects.

\subsection{Interpretation}
Explain what the result means technically. Distinguish measured fact from engineering inference.

\subsection{Conclusion}
State whether the hypothesis was confirmed, rejected, or left unresolved.

\subsection{Next action}
Specify the follow-up experiment, design change, additional instrumentation, or decision triggered by this run.

\subsection{Evidence attachments}
\begin{RDChecklist}
\item Commit references.
\item Benchmark exports.
\item Logs or screenshots.
\item Diagram revision references.
\item PR, ticket, or release references.
\end{RDChecklist}

\clearpage
\section{Repeatable blank run pages}
\subsection{Blank experiment run page}
\RDCallout{Use}{Duplicate this section for each independent run or move it into a separate run-specific file if a project generates heavy volume.}

Run ID: \rule{5cm}{0.4pt}\\[0.15cm]
Date or range: \rule{5cm}{0.4pt}\\[0.15cm]
Engineer(s): \rule{8cm}{0.4pt}\\[0.15cm]
Project / uncertainty ref: \rule{8cm}{0.4pt}

\subsubsection{Technical objective}
\vspace{2.3cm}

\subsubsection{Background}
\vspace{3.0cm}

\subsubsection{Hypothesis}
\vspace{1.8cm}

\subsubsection{Method}
\vspace{4.5cm}

\subsubsection{Results and observations}
\vspace{5.5cm}

\subsubsection{Conclusion and next action}
\vspace{3.2cm}

\RDEndMatter
\end{document}
